% 图契约：本题的源、边界和推导逐项对应；箭头符号按图注解释。
\documentclass[tikz,border=7pt]{standalone}
\usepackage{fontspec,xeCJK,amsmath,amssymb}
\IfFontExistsTF{Times New Roman}{\setmainfont{Times New Roman}}{\setmainfont{TeX Gyre Termes}}
\IfFontExistsTF{Heiti SC}{\setCJKmainfont[AutoFakeBold=1.5]{Heiti SC}}{\setCJKmainfont[AutoFakeBold=1.5]{FandolHei-Regular}}
\usetikzlibrary{arrows.meta,bending,calc,patterns,positioning}
\definecolor{ink}{HTML}{243745}
\definecolor{blue}{HTML}{275DAB}
\definecolor{green}{HTML}{177C63}
\definecolor{red}{HTML}{BD4935}
\definecolor{orange}{HTML}{AA6C16}
\definecolor{purple}{HTML}{715393}
\tikzset{every node/.style={font=\fontsize{9.3}{12.5}\selectfont,text=ink},
axis/.style={-{Stealth[length=4pt]},line width=.65pt,draw=ink!65},
vec/.style={-{Stealth[length=5pt]},line width=1pt,draw=blue},
short/.style={-{Stealth[length=3pt]},line width=.8pt},
guide/.style={dashed,draw=ink!40,line width=.5pt},
card/.style={rounded corners=2pt,draw=blue!65,fill=blue!4,align=center,text width=6.5cm,inner xsep=9pt,inner ysep=7pt},
title/.style={font=\fontsize{10.5}{14}\selectfont\bfseries,text width=15cm,align=center},
label/.style={font=\fontsize{8.4}{11}\selectfont,align=center},
note/.style={label,text width=6.2cm}}
\begin{document}
\begin{tikzpicture}[x=1cm,y=1cm]
\path[use as bounding box] (-8,-3.5) rectangle(8,3.5);
\node[title] at (0,3.15) {M2-156\quad 平方余弦拆成零阶与二阶};
\begin{scope}[shift={(-4,0)}]
\begin{scope}\clip (0,0) circle(1.6);\fill[blue!7] (-1.7,-1.7) rectangle(1.7,1.7);\fill[red!12] (0,0)--(35.264:1.7) arc(35.264:144.736:1.7)--cycle;\fill[red!12] (0,0)--(215.264:1.7) arc(215.264:324.736:1.7)--cycle;\end{scope}\draw[ink!55] (0,0)circle(1.6);\draw[axis] (0,-1.95)--(0,2.1) node[above]{$z$};\draw[guide](0,0)--(35:1.6) node[midway,above left]{$R$};\node[label] at(0,-2.3){红、蓝区示意该角模的正、负部分};
\node[note] at (0,-2.85) {着色仅表示二阶部分，完整边界还要加上平均值。};
\end{scope}
\node[card,draw=blue!65,fill=blue!4] at (3.9,1.8) {$\cos^2\theta=\frac13+\frac23P_2(\cos\theta)$};
\node[card,draw=green!65,fill=green!4] at (3.9,0) {$\displaystyle\varphi=V_0\left[\frac13+\frac23(r/R)^2P_2\right]$};
\node[card,draw=orange!65,fill=orange!4] at (3.9,-1.8) {零阶为常量，二阶配 $r^2$；\\球心电势为 $V_0/3$。};
\end{tikzpicture}\end{document}